\clearpage
\fancyhead[L]{\small COMET.}
\hypertarget{native-455-comet}{}
\begin{center}
{\large COMET.}
\end{center}

\begin{center}
\rule{6cm}{0.4pt}\\
Astronomische Nachrichten. Band III. Nr. 59. Seite 179..182. 1824 Juni.\\
\rule{6cm}{0.4pt}
\end{center}

\begin{center}
Beobachtungen des Cometen von 1824 in Göttingen.
\end{center}

Diese Beobachtungen sind zwar alle am Meridiankreise gemacht, bei den obern
Culminationen; allein die grosse Lichtschwäche des Cometen machte es
unmöglich, ihn bei der geringsten Fadenbeleuchtung zu sehen, so dass
wechselweise die Fäden beim Oeffnen und der Comet beim Schliessen der
Beleuchtung gesehen werden mussten. Dieser Umstand ist Schuld, dass die
Beobachtungen lange nicht so genau sein können, als sie unter günstigern
Verhältnissen gewesen sein würden. Die Beobachtung aller frühern
Culminationen, obgleich ich jedesmal mich dazu in Bereitschaft gesetzt hatte,
wurden vom Wetter vereitelt. Ich bemerke noch, dass nicht immer Antritte an
allen 7 Fäden beobachtet oder geschätzt sind, und dass daher die angesetzte
Zeit, welche immer dem Mittel der beobachteten Antritte entspricht, und auf
welche die Declinationen, falls sie nicht gleichzeitig eingestellt waren,
reducirt worden sind, nicht immer mit den Durchgängen durch den Meridian
selbst übereinstimmen.

\begin{center}
\begingroup\small\setlength{\tabcolsep}{4pt}
\begin{tabular}{lr|c|c|c}
\multicolumn{2}{c|}{} & \shortstack{Göttingen\\Mittlere Zeit} & \shortstack{Gerade\\Aufsteigung} & \shortstack{Abweichung\\Nordlich.} \\
\hline
Januar & 31 & $13^{\rm u}\,28^{\rm m}\,20^{\rm s}5$ & $152^{\circ}\,16'\,29''1$ & $72^{\circ}\,38'\,35''3$ \\
Februar & 3 & $11\quad35\quad7.1$ & $131\quad23\quad10$ & $69\quad14\quad8.6$ \\
 & 7 & $10\quad38\quad2.8$ & $116\quad12\quad49.6$ & $63\quad14\quad55.9$ \\
 & 8 & $10\quad24\quad5.3$ & $113\quad50\quad49.2$ & $61\quad46\quad15.0$ \\
 & 19 & $8\quad52\quad23.5$ & $101\quad49\quad31.2$ & $48\quad48\quad51.1$ \\
 & 20 & $8\quad48$ &  & $47\quad55$ \quad :\,: \\
 & 27 & $8\quad11\quad46.1$ & .\,$99\quad26\quad50.5$ & $42\quad48\quad49.0$ \\
 & 28 & $8\quad6\quad47.2$ & $99\quad18\quad7.7$ & $42\quad12\quad46.2$ \\
März & 2 & $7\quad53\quad59.0$ & $99\quad0\quad39.4$ & $40\quad32\quad2.5$ \\
\end{tabular}
\endgroup
\end{center}

\clearpage
\fancyhead[L]{\small GAUSS AN BODE.}
\SourceLeaf{456}

\begin{center}
{\large GAUSS AN BODE.}
\end{center}

\begin{center}
\rule{6cm}{0.4pt}\\
Astronomisches Jahrbuch für 1828. Seite 171..172. Berlin 1825.\\
\rule{6cm}{0.4pt}
\end{center}

\begin{flushright}
Göttingen, den 5. August 1825.
\end{flushright}

Diese Beobachtungen gewinnen vielleicht dadurch an Interesse, dass diese
Planeten 1825 zum erstenmale wieder sehr nahe in derselben Gegend waren, wo
sie 1802 resp. entdeckt und wieder aufgefunden wurden, auch ist vielleicht
wenigstens die Pallas sonst nirgends beobachtet, da die bekannt gemachte
Ephemeride so sehr viel abweicht: mir selbst hatte diesmal die Zeit gefehlt,
nach meiner Theorie eine solche zu berechnen.

\begin{center}
\emph{Meridian-Beobachtungen der Pallas auf der Sternwarte in Göttingen.}
\end{center}

\begin{center}
\begingroup\small\setlength{\tabcolsep}{4pt}
\begin{tabular}{lr|c|c|c}
\multicolumn{2}{c|}{1825} & \shortstack{Mittlere Zeit\\in Göttingen} & \shortstack{Gerade\\Aufsteigung} & Abweichung \\
\hline
März & 3 & $13^{\rm u}\,4^{\rm m}\,23^{\rm s}7$ & $177^{\circ}\,34'\,24''3$ & $0^{\circ}\,24'\,4''1$ Südl. \\
 & 4 & $12\quad59\quad49.6$ & $177\quad24\quad50.7$ & $0\quad0\quad18.2$ Nordl. \\
 & 7 & $12\quad46\quad4.0$ & $176\quad55\quad18.1$ & $1\quad14\quad7.2$ \\
 & 8 & $12\quad41\quad27.9$ & $176\quad45\quad12.3$ & $1\quad38\quad52.6$ \\
 & 9 & $12\quad36\quad51.3$ & $176\quad35\quad0.3$ & $2\quad3\quad40.3$ \\
 & 10 & $12\quad32\quad14.3$ & $176\quad24\quad43.2$ & $2\quad28\quad29.3$ \\
 & 18 & $11\quad55\quad14.1$ & $175\quad1\quad15.0$ & $5\quad44\quad59.3$ \\
 & 19 & $11\quad50\quad37.3$ & $174\quad50\quad59.1$ & $6\quad8\quad56.3$ \\
 & 20 & $11\quad46\quad0.6$ & $174\quad40\quad46.2$ & $6\quad32\quad41.4$ \\
 & 25 & $11\quad23\quad6.1$ & $173\quad51\quad53.7$ & $8\quad27\quad50.6$ \\
 & 27 & $11\quad14\quad1.1$ & $173\quad33\quad32.8$ & $9\quad11\quad51.7$ \\
 & 29 & $11\quad4\quad59.8$ & $173\quad16\quad8.2$ & $9\quad54\quad29.8$ \\
April & 6 & $10\quad29\quad36.6$ & $172\quad16\quad59.2$ & $12\quad29\quad56.1$ \\
 & 7 & $10\quad25\quad16.7$ & $172\quad10\quad57.4$ & $12\quad47\quad32.9$ \\
\end{tabular}
\endgroup
\end{center}

\begin{center}
\emph{Meridian-Beobachtungen der Ceres auf der Sternwarte in Göttingen.}
\end{center}

\begin{center}
\begingroup\small\setlength{\tabcolsep}{4pt}
\begin{tabular}{lr|c|c|c}
\multicolumn{2}{c|}{1825} & \shortstack{Mittlere Zeit\\in Göttingen} & \shortstack{Gerade\\Aufsteigung} & Abweichung \\
\hline
März & 9 & $13^{\rm u}\,1^{\rm m}\,7^{\rm s}5$ & $182^{\circ}\,40'\,4''6$ & $17^{\circ}\,35'\,18''5$ Nordl. \\
 & 10 & $12\quad56\quad23.3$ & $182\quad27\quad57.4$ & $17\quad41\quad20.9$ \\
 & 18 & $12\quad18\quad11.1$ & $180\quad46\quad26.4$ & $18\quad24\quad4.4$ \\
 & 19 & $12\quad13\quad23.5$ & $180\quad33\quad28.9$ & $18\quad28\quad34.0$ \\
 & 20 & $12\quad8\quad36.8$ & $180\quad20\quad32.2$ & $18\quad32\quad50.0$ \\
 & 29 & $11\quad25\quad38.4$ & $178\quad26\quad37.3$ & $19\quad1\quad18.0$ \\
April & 7 & $10\quad43\quad31.3$ & $176\quad45\quad21.7$ & $19\quad10\quad5.9$ \\
\end{tabular}
\endgroup
\end{center}
